\chapter{{\fontsize{14pt}{21pt}\selectfont\MakeUppercase{Текущее состояние процесса (как решается сейчас)}}}
\label{ch:current-state}

Текущее состояние процесса управления цифровым присутствием \abbr{малого и среднего бизнеса}{МСБ}~--- модель AS-IS~--- характеризуется следующими особенностями.

\section{Как проблема целевой аудитории решается в настоящее время}

\begin{enumerate}
  \item Владелец или сотрудник бизнеса вручную авторизуется на каждой платформе (Яндекс.Бизнес, 2ГИС, VK, Telegram) через отдельный веб-интерфейс.
  \item При изменении бизнес-информации (часы работы, адрес, контакты) обновление выполняется последовательно на каждой площадке с~индивидуальными формами редактирования.
  \item Мониторинг отзывов и сообщений требует поочерёдной проверки каждой платформы; единого потока событий не существует.
  \item Отсутствует <<эталонный>> источник данных~--- актуальная информация хранится только на самих площадках, что затрудняет обнаружение расхождений.
  \item Результаты действий не журналируются централизованно~--- невозможно отследить, кто и когда обновил данные на конкретной платформе.
\end{enumerate}

Укрупнённая схема текущего процесса (AS-IS) может быть сгенерирована из исходника \texttt{as-is.mmd} (Mermaid) в~директории пояснительной записки ВКР; в~данном отчёте ограничиваемся текстовым описанием. Основные недостатки~--- высокая трудоёмкость ручных операций, риск рассогласования данных между площадками и отсутствие централизованного контроля~--- формируют предпосылки к~проектированию автоматизированной системы.
